\documentclass[aspectratio=169]{beamer}
\usepackage[utf8]{inputenc}
\usepackage{amsmath,amssymb}
\usepackage{graphicx}
\usepackage{minted}
\usetheme{Madrid}

\title{Chapter 2: Supervised Learning - Linear Models}
\author{Dr. Duc-Quang VU}
\institute{Machine Learning Course}
\date{\today}

\begin{document}

\frame{\titlepage}

\begin{frame}
\frametitle{Outline}
\tableofcontents
\end{frame}

\section{2.1 Linear Regression}

\begin{frame}{Linear Regression: Introduction (1)}
\begin{itemize}
    \item Supervised learning task for predicting a continuous numerical output.
    \item Goal: Establish a relationship model between input features ($X$) and a continuous target variable ($y$).
    \item Essentially, we are fitting a line (or a hyperplane in higher dimensions) to a set of data points.
\end{itemize}
\end{frame}

\begin{frame}{Linear Regression: Introduction (2) - Use Cases}
\begin{itemize}
    \item \textbf{Real Estate:} Predicting house prices based on square footage, number of bedrooms, and location.
    \item \textbf{Finance:} Forecasting company revenue or stock prices based on historical trends.
    \item \textbf{Healthcare:} Estimating a patient's blood pressure based on age, weight, and BMI.
    \item \textbf{Economics:} Predicting GDP growth based on various economic indicators.
\end{itemize}
\end{frame}

\begin{frame}{Mathematical Notation \& Setup}
\begin{itemize}
    \item $m$: Number of training examples (size of the dataset).
    \item $n$: Number of features (input variables).
    \item $x$: Input variables / features. $x \in \mathbb{R}^n$.
    \item $y$: Output variable / target variable. $y \in \mathbb{R}$.
    \item $(x^{(i)}, y^{(i)})$ represents the $i^{th}$ training example.
    \item $x_j^{(i)}$ is the value of feature $j$ in the $i^{th}$ training example.
\end{itemize}
\end{frame}

\begin{frame}{The Hypothesis Function (1)}
\begin{itemize}
    \item In supervised learning, the model is often called a "hypothesis", denoted as $h$.
    \item For univariate linear regression (one feature), the hypothesis is a straight line equation:
    $$ h_\theta(x) = \theta_0 + \theta_1 x $$
    \item $\theta_0$: y-intercept (bias).
    \item $\theta_1$: slope (weight of feature $x$).
    \item $\theta_i$ are called the \textbf{parameters} or \textbf{weights} of the model.
\end{itemize}
\end{frame}

\begin{frame}{The Hypothesis Function (2) - Multivariate}
\begin{itemize}
    \item For multiple linear regression (multiple features $x_1, x_2, \dots, x_n$):
    $$ h_\theta(x) = \theta_0 + \theta_1 x_1 + \theta_2 x_2 + \dots + \theta_n x_n $$
    \item For convenience, we define $x_0 = 1$. Then:
    $$ x = \begin{bmatrix} x_0 \\ x_1 \\ \vdots \\ x_n \end{bmatrix} \quad \text{and} \quad \theta = \begin{bmatrix} \theta_0 \\ \theta_1 \\ \vdots \\ \theta_n \end{bmatrix} $$
\end{itemize}
\end{frame}

\begin{frame}{The Hypothesis Function (3) - Vectorization}
\begin{itemize}
    \item Using linear algebra, we can write the hypothesis compactly as an inner product (dot product):
    $$ h_\theta(x) = \theta_0 x_0 + \theta_1 x_1 + \dots + \theta_n x_n = \theta^T x $$
    \item Where $\theta^T$ is the transpose of the parameter vector $\theta$.
    \item Vectorization is crucial for efficient computation in Machine Learning (e.g., using NumPy in Python).
\end{itemize}
\end{frame}

\begin{frame}{Defining "Goodness of Fit": The Cost Function (1)}
\begin{itemize}
    \item How do we choose the best parameters $\theta$?
    \item We need a metric to evaluate how well our hypothesis $h_\theta(x)$ fits the actual data $y$.
    \item We define a \textbf{Cost Function} (or Loss Function), $J(\theta)$, which measures the difference between our model's predictions and the true values.
\end{itemize}
\end{frame}

\begin{frame}{The Cost Function (2) - Mean Squared Error}
\begin{itemize}
    \item The most common cost function for linear regression is the \textbf{Mean Squared Error (MSE)}.
    \item Formula:
    $$ J(\theta) = \frac{1}{2m} \sum_{i=1}^{m} (h_\theta(x^{(i)}) - y^{(i)})^2 $$
    \item Why squared?
        \begin{itemize}
            \item Prevents positive and negative errors from canceling out.
            \item Penalizes larger errors more heavily than smaller ones.
        \end{itemize}
    \item Why $\frac{1}{2m}$? The factor of $2$ makes the derivative cleaner later on.
\end{itemize}
\end{frame}

\begin{frame}{The Optimization Objective}
\begin{itemize}
    \item Our primary goal in Linear Regression is to find the set of parameters $\theta$ that yields the smallest possible error.
    \item Mathematically, this is an optimization problem:
    $$ \min_\theta J(\theta) $$
    \item The $\theta$ that minimizes $J(\theta)$ will give us the line of best fit for our training data.
\end{itemize}
\end{frame}

\begin{frame}{Visualizing the Cost Function (1D)}
\begin{itemize}
    \item Consider simplified hypothesis $h_\theta(x) = \theta_1 x$ (setting $\theta_0 = 0$).
    \item The cost function $J(\theta_1)$ is a function of a single variable.
    \item If we plot $J(\theta_1)$ against $\theta_1$, we get a parabola (a convex shape).
    \item The minimum of this parabola corresponds to the optimal value of $\theta_1$.
\end{itemize}
\end{frame}

\begin{frame}{Visualizing the Cost Function (2D)}
\begin{itemize}
    \item With both $\theta_0$ and $\theta_1$, the cost function $J(\theta_0, \theta_1)$ forms a 3D surface.
    \item It looks like a bowl (it is a \textbf{convex function}).
    \item The lowest point of the bowl represents the global minimum of the cost function.
    \item Contour plots (2D slices of the 3D surface) are often used to visualize this. Concentric ellipses represent equal cost.
\end{itemize}
\end{frame}

\begin{frame}{Finding the Minimum: Optimization Algorithms}
\begin{itemize}
    \item We have defined the objective: $\min_\theta J(\theta)$. How do we actually compute the optimal $\theta$?
    \item Two primary approaches for Linear Regression:
    \begin{enumerate}
        \item \textbf{Iterative Method:} Gradient Descent. (Applicable to many ML algorithms).
        \item \textbf{Analytical Method:} Normal Equation. (Specific to Linear Regression).
    \end{enumerate}
\end{itemize}
\end{frame}

\begin{frame}{Gradient Descent: Intuition}
\begin{itemize}
    \item Imagine you are blindfolded at the top of a hill and want to reach the bottom (the minimum).
    \item Strategy:
    \begin{enumerate}
        \item Feel the slope of the ground around you (calculate the gradient/derivative).
        \item Take a step in the direction of the steepest descent (downhill).
        \item Repeat until you reach a flat area (minimum).
    \end{enumerate}
\end{itemize}
\end{frame}

\begin{frame}{Gradient Descent: The Algorithm}
\begin{itemize}
    \item Start with initial guesses for $\theta$ (e.g., all zeros).
    \item Repeat until convergence:
    $$ \theta_j := \theta_j - \alpha \frac{\partial}{\partial \theta_j} J(\theta_0, \dots, \theta_n) \quad (\text{for } j=0, 1, \dots, n) $$
    \item \textbf{Crucial details:}
        \begin{itemize}
            \item $:=$ denotes assignment.
            \item $\alpha$ is the \textbf{learning rate} (step size).
            \item $\frac{\partial}{\partial \theta_j} J(\theta)$ is the partial derivative of the cost function with respect to parameter $\theta_j$.
            \item \textbf{Simultaneous update:} You must compute new values for all $\theta_j$ before updating them.
        \end{itemize}
\end{itemize}
\end{frame}

\begin{frame}{Gradient Descent for Linear Regression}
\begin{itemize}
    \item We need to compute the partial derivative term for the MSE cost function:
    $$ \frac{\partial}{\partial \theta_j} J(\theta) = \frac{\partial}{\partial \theta_j} \left( \frac{1}{2m} \sum_{i=1}^{m} (h_\theta(x^{(i)}) - y^{(i)})^2 \right) $$
    \item Applying the chain rule, we get:
    $$ \frac{\partial}{\partial \theta_j} J(\theta) = \frac{1}{m} \sum_{i=1}^{m} (h_\theta(x^{(i)}) - y^{(i)}) \cdot x_j^{(i)} $$
\end{itemize}
\end{frame}

\begin{frame}{The Gradient Descent Update Rules}
\begin{itemize}
    \item Substituting the derivative back into the algorithm:
    \item Repeat until convergence:
        \begin{align*}
        \theta_0 &:= \theta_0 - \alpha \frac{1}{m} \sum_{i=1}^{m} (h_\theta(x^{(i)}) - y^{(i)}) \cdot 1 \\
        \theta_1 &:= \theta_1 - \alpha \frac{1}{m} \sum_{i=1}^{m} (h_\theta(x^{(i)}) - y^{(i)}) \cdot x_1^{(i)} \\
        &\vdots \\
        \theta_j &:= \theta_j - \alpha \frac{1}{m} \sum_{i=1}^{m} (h_\theta(x^{(i)}) - y^{(i)}) \cdot x_j^{(i)}
        \end{align*}
\end{itemize}
\end{frame}

\begin{frame}{The Learning Rate $\alpha$}
\begin{itemize}
    \item The learning rate $\alpha$ determines the size of the steps we take downhill.
    \item \textbf{If $\alpha$ is too small:} Gradient descent can be very slow, taking many iterations to converge.
    \item \textbf{If $\alpha$ is too large:} Gradient descent can overshoot the minimum. It may fail to converge, or even diverge (error increases).
    \item How to choose? Try values on a log scale (e.g., $0.001, 0.003, 0.01, 0.03, 0.1, 0.3$) and plot $J(\theta)$ over iterations.
\end{itemize}
\end{frame}

\begin{frame}{Convexity of MSE}
\begin{itemize}
    \item The MSE cost function for linear regression is a \textbf{convex function}.
    \item This means it is bowl-shaped and has no local optima, only one global optimum.
    \item Therefore, gradient descent on MSE is guaranteed to converge to the global minimum (as long as $\alpha$ is not too large).
\end{itemize}
\end{frame}

\begin{frame}{Feature Scaling: The Problem}
\begin{itemize}
    \item Consider two features: $x_1$ = size in sq ft (range: 1000-2000), $x_2$ = number of bedrooms (range: 1-5).
    \item The contours of the cost function $J(\theta)$ will be very tall and skinny ellipses.
    \item Gradient descent will oscillate wildly back and forth before eventually finding its way to the minimum. It takes a very long time.
\end{itemize}
\end{frame}

\begin{frame}{Feature Scaling: The Solution}
\begin{itemize}
    \item We scale the features so they are on roughly the same scale, typically $-1 \le x_i \le 1$.
    \item This makes the contours of the cost function more circular, allowing gradient descent to converge much faster via a more direct path to the minimum.
\end{itemize}
\end{frame}

\begin{frame}{Feature Scaling: Methods}
\begin{itemize}
    \item \textbf{Mean Normalization:}
    $$ x_i := \frac{x_i - \mu_i}{s_i} $$
    Where $\mu_i$ is the average value of feature $i$ in training set, and $s_i$ is the range (max - min).
    \item \textbf{Standardization (Z-score):}
    $$ x_i := \frac{x_i - \mu_i}{\sigma_i} $$
    Where $\sigma_i$ is the standard deviation. (Most common in practice).
\end{itemize}
\end{frame}

\begin{frame}{The Normal Equation (1)}
\begin{itemize}
    \item For Linear Regression, we can solve for the optimal $\theta$ analytically using calculus (setting derivatives to zero).
    \item This gives us the \textbf{Normal Equation}:
    $$ \theta = (X^T X)^{-1} X^T y $$
    \item Where:
    \begin{itemize}
        \item $X$ is the \textbf{Design Matrix} (size $m \times (n+1)$). Each row is a training example $(x^{(i)})^T$.
        \item $y$ is the vector of target values (size $m \times 1$).
    \end{itemize}
\end{itemize}
\end{frame}

\begin{frame}{Design Matrix $X$ Example}
\begin{itemize}
    \item Suppose we have 3 training examples and 2 features ($x_1, x_2$).
    $$ X = \begin{bmatrix}
    1 & x_1^{(1)} & x_2^{(1)} \\
    1 & x_1^{(2)} & x_2^{(2)} \\
    1 & x_1^{(3)} & x_2^{(3)}
    \end{bmatrix}, \quad
    y = \begin{bmatrix}
    y^{(1)} \\ y^{(2)} \\ y^{(3)}
    \end{bmatrix} $$
    \item The first column of $1$s corresponds to the intercept term $x_0 = 1$.
\end{itemize}
\end{frame}

\begin{frame}{Gradient Descent vs. Normal Equation}
\begin{table}
\centering
\small
\begin{tabular}{|p{0.45\textwidth}|p{0.45\textwidth}|}
\hline
\textbf{Gradient Descent} & \textbf{Normal Equation} \\
\hline
Need to choose learning rate $\alpha$ & No need to choose $\alpha$ \\
Needs many iterations & No iterations (compute once) \\
Works well even when $n$ is large & Slow if $n$ is very large ($O(n^3)$ for inversion) \\
Requires feature scaling & Doesn't require feature scaling \\
Applicable to many algorithms & Only applies to Linear Regression \\
\hline
\end{tabular}
\end{table}
\end{frame}

\begin{frame}{Polynomial Regression (1)}
\begin{itemize}
    \item What if a straight line doesn't fit the data well (e.g., the relationship is curved)?
    \item We can use \textbf{Polynomial Regression}, which is a special case of multiple linear regression.
    \item We create new features by taking powers of existing features.
\end{itemize}
\end{frame}

\begin{frame}{Polynomial Regression (2)}
\begin{itemize}
    \item Example: Original feature $x$ (size of house).
    \item Hypothesis: $h_\theta(x) = \theta_0 + \theta_1 x + \theta_2 x^2 + \theta_3 x^3$.
    \item We treat this as multiple linear regression with features:
    \begin{itemize}
        \item $x_1 = x$
        \item $x_2 = x^2$
        \item $x_3 = x^3$
    \end{itemize}
    \item Important: Feature scaling becomes absolutely critical here, as ranges of $x, x^2, x^3$ will vary drastically.
\end{itemize}
\end{frame}

\begin{frame}{Evaluating Model Performance: $R^2$ Score}
\begin{itemize}
    \item How do we know if our model is good after training?
    \item The \textbf{Coefficient of Determination ($R^2$ score)} is a common metric.
    \item $R^2 = 1 - \frac{\text{Sum of Squared Residuals (SSR)}}{\text{Total Sum of Squares (SST)}}$
    \item Interpretation: It represents the proportion of variance in the dependent variable ($y$) that is predictable from the independent variables ($X$).
    \item $R^2 \approx 1$: Perfect fit. $R^2 \approx 0$: Model is no better than predicting the mean of $y$.
\end{itemize}
\end{frame}

\begin{frame}[fragile]{Linear Regression: Python Demo Setup}
\begin{itemize}
    \item We will use \texttt{scikit-learn} (sklearn), a standard machine learning library in Python.
    \item It abstracts away the math (Gradient Descent/Normal Equation) into simple API calls.
\end{itemize}
\begin{minted}[fontsize=\footnotesize]{python}
import numpy as np
import matplotlib.pyplot as plt
from sklearn.linear_model import LinearRegression
from sklearn.model_selection import train_test_split
from sklearn.metrics import mean_squared_error, r2_score
\end{minted}
\end{frame}

\begin{frame}[fragile]{Linear Regression: Python Demo Code}
\begin{minted}[fontsize=\scriptsize]{python}
# 1. Generate dummy data with some noise
np.random.seed(42)
X = 2 * np.random.rand(100, 1) # Feature
y = 4 + 3 * X + np.random.randn(100, 1) # Target line: y = 4 + 3x

# 2. Split data into training and testing sets
X_train, X_test, y_train, y_test = train_test_split(
    X, y, test_size=0.2, random_state=42
)

# 3. Initialize and train the model
model = LinearRegression()
model.fit(X_train, y_train) # Model learns parameters

# 4. Make predictions on test set
y_pred = model.predict(X_test)
\end{minted}
\end{frame}

\begin{frame}[fragile]{Linear Regression: Python Demo Results}
\begin{minted}[fontsize=\footnotesize]{python}
# 5. Evaluate the model
mse = mean_squared_error(y_test, y_pred)
r2 = r2_score(y_test, y_pred)

print(f"Mean Squared Error: {mse:.4f}")
print(f"R-squared Score: {r2:.4f}")

# Model parameters
print(f"Intercept (theta_0): {model.intercept_[0]:.4f}")
print(f"Coefficient (theta_1): {model.coef_[0][0]:.4f}")
# True parameters were 4 and 3. Model should be close.
\end{minted}
\end{frame}

\begin{frame}{Linear Regression: Summary (2.1)}
    \begin{itemize}
        \item A foundational algorithm for predicting continuous values.
        \item Models the relationship as a linear combination of features.
        \item \textbf{Cost Function:} Mean Squared Error (MSE).
        \item \textbf{Optimization:} Gradient Descent (iterative) or Normal Equation (analytical).
        \item \textbf{Extensions:} Can handle non-linear relationships via Polynomial Regression.
        \item \textbf{Best Practices:} Always scale features when using Gradient Descent or Polynomial features.
    \end{itemize}
\end{frame}


\section{2.2 Logistic Regression}

\begin{frame}{Logistic Regression: Introduction (1)}
\begin{itemize}
    \item Despite having "regression" in its name, it is a \textbf{classification} algorithm.
    \item Used when the target variable $y$ is discrete (categorical) rather than continuous.
    \item Focuses primarily on \textbf{Binary Classification}: predicting one of two possible outcomes.
    \item The target variable $y \in \{0, 1\}$.
        \begin{itemize}
            \item $0$: Negative Class (e.g., Benign tumor, Not spam).
            \item $1$: Positive Class (e.g., Malignant tumor, Spam).
        \end{itemize}
\end{itemize}
\end{frame}

\begin{frame}{Logistic Regression: Introduction (2) - Use Cases}
\begin{itemize}
    \item \textbf{Email Filtering:} Classifying an incoming email as Spam (1) or Not Spam (0).
    \item \textbf{Medical Diagnosis:} Determining if a patient has a specific disease based on symptoms.
    \item \textbf{Banking/Finance:} Predicting if a customer will default on a loan (Credit Scoring) or detecting fraudulent transactions.
    \item \textbf{Marketing:} Predicting if a user will click on an ad (Click-Through Rate prediction).
\end{itemize}
\end{frame}

\begin{frame}{Why Not Linear Regression for Classification?}
\begin{itemize}
    \item Applying linear regression to classification tasks is problematic.
    \item \textbf{Issue 1: Output Range.} Linear regression outputs values $h_\theta(x) \in (-\infty, \infty)$, but we want probabilities between $0$ and $1$.
    \item \textbf{Issue 2: Sensitivity to Outliers.} Adding a single data point far to the right can drastically shift the regression line, changing the decision threshold and causing misclassifications of previously correct examples.
    \item We need a hypothesis function that is bounded: $0 \le h_\theta(x) \le 1$.
\end{itemize}
\end{frame}

\begin{frame}{Logistic Regression: Hypothesis Representation}
\begin{itemize}
    \item To bound our output between 0 and 1, we apply a mathematical function called the \textbf{Logistic Function} (or Sigmoid function) to the linear combination of features.
    \item Linear model: $z = \theta^T x$
    \item Logistic Regression Hypothesis: 
    $$ h_\theta(x) = g(\theta^T x) $$
    \item Where $g(z)$ is the Sigmoid function:
    $$ g(z) = \frac{1}{1 + e^{-z}} $$
\end{itemize}
\end{frame}

\begin{frame}{The Sigmoid Function}
\begin{itemize}
    \item Formula: $g(z) = \frac{1}{1 + e^{-z}}$
    \item Key Properties:
    \begin{itemize}
        \item As $z \rightarrow +\infty$, $e^{-z} \rightarrow 0$, so $g(z) \rightarrow 1$.
        \item As $z \rightarrow -\infty$, $e^{-z} \rightarrow \infty$, so $g(z) \rightarrow 0$.
        \item When $z = 0$, $e^{0} = 1$, so $g(0) = 0.5$.
    \end{itemize}
    \item The Sigmoid function smoothly maps any real-valued number into the range $(0, 1)$, creating an 'S' shaped curve.
\end{itemize}
\end{frame}

\begin{frame}{Probabilistic Interpretation}
\begin{itemize}
    \item We interpret the output of the hypothesis $h_\theta(x)$ as the \textbf{estimated probability} that $y=1$ on input $x$.
    \item Notation:
    $$ h_\theta(x) = P(y=1 | x; \theta) $$
    (Probability that $y=1$, given $x$, parameterized by $\theta$)
    \item Since $y$ must be either 0 or 1, the probabilities must sum to 1:
    $$ P(y=0 | x; \theta) = 1 - P(y=1 | x; \theta) = 1 - h_\theta(x) $$
\end{itemize}
\end{frame}

\begin{frame}{Decision Boundary (1)}
\begin{itemize}
    \item To translate the probability $h_\theta(x)$ into a discrete class prediction (0 or 1), we use a threshold.
    \item Standard Threshold = $0.5$:
    \begin{itemize}
        \item Predict $y=1$ if $h_\theta(x) \ge 0.5$.
        \item Predict $y=0$ if $h_\theta(x) < 0.5$.
    \end{itemize}
    \item Looking at the Sigmoid function, $g(z) \ge 0.5$ whenever $z \ge 0$.
    \item Therefore: Predict $y=1$ if $\theta^T x \ge 0$.
\end{itemize}
\end{frame}

\begin{frame}{Decision Boundary (2)}
\begin{itemize}
    \item The equation $\theta^T x = 0$ defines the \textbf{Decision Boundary}.
    \item It is the line (or hyperplane in higher dimensions) that separates the region where the model predicts $y=1$ from the region where it predicts $y=0$.
    \item The decision boundary is a property of the hypothesis and parameters $\theta$, not the dataset itself.
\end{itemize}
\end{frame}

\begin{frame}{Non-linear Decision Boundaries}
\begin{itemize}
    \item Just like in Linear Regression, we can create more complex, non-linear decision boundaries by adding polynomial features.
    \item Example: Let features be $x_1, x_2, x_1^2, x_2^2$.
    \item Hypothesis: $h_\theta(x) = g(\theta_0 + \theta_1 x_1 + \theta_2 x_2 + \theta_3 x_1^2 + \theta_4 x_2^2)$
    \item The decision boundary ($\theta^T x = 0$) could now be a circle or an ellipse (e.g., $x_1^2 + x_2^2 = 1$).
\end{itemize}
\end{frame}

\begin{frame}{The Cost Function: Why not MSE?}
\begin{itemize}
    \item If we use the Mean Squared Error (MSE) cost function from Linear Regression for Logistic Regression, we run into a major problem.
    \item $J(\theta) = \frac{1}{2m} \sum_{i=1}^{m} (g(\theta^T x^{(i)}) - y^{(i)})^2$
    \item Because the Sigmoid function $g(z)$ is highly non-linear, this MSE cost function becomes \textbf{non-convex}.
    \item It will have many local minima, and Gradient Descent is not guaranteed to find the global minimum.
\end{itemize}
\end{frame}

\begin{frame}{The Logistic Regression Cost Function (Cross-Entropy)}
\begin{itemize}
    \item We need a cost function that is convex. We define the cost for a single training example as:
    $$ Cost(h_\theta(x), y) = \begin{cases} -\log(h_\theta(x)) & \text{if } y = 1 \\ -\log(1 - h_\theta(x)) & \text{if } y = 0 \end{cases} $$
    \item Also known as \textbf{Log Loss} or Binary Cross-Entropy.
\end{itemize}
\end{frame}

\begin{frame}{Intuition behind Cross-Entropy (1)}
\begin{itemize}
    \item Case 1: True label $y = 1$.
        \begin{itemize}
            \item Cost is $-\log(h_\theta(x))$.
            \item If the model predicts $h_\theta(x) = 1$ (perfect!), cost $= -\log(1) = 0$.
            \item If the model predicts $h_\theta(x) \rightarrow 0$ (completely wrong), cost $\rightarrow \infty$.
            \item It heavily penalizes the model for being very confident and very wrong.
        \end{itemize}
\end{itemize}
\end{frame}

\begin{frame}{Intuition behind Cross-Entropy (2)}
\begin{itemize}
    \item Case 2: True label $y = 0$.
        \begin{itemize}
            \item Cost is $-\log(1 - h_\theta(x))$.
            \item If the model predicts $h_\theta(x) = 0$ (perfect!), cost $= -\log(1) = 0$.
            \item If the model predicts $h_\theta(x) \rightarrow 1$ (completely wrong), cost $\rightarrow \infty$.
        \end{itemize}
\end{itemize}
\end{frame}

\begin{frame}{Simplified Cost Function Equation}
\begin{itemize}
    \item We can compress the two cases into a single, elegant equation (since $y$ is either 0 or 1, one term always cancels out):
    $$ Cost(h_\theta(x), y) = -y \log(h_\theta(x)) - (1 - y) \log(1 - h_\theta(x)) $$
    \item The total cost function over the entire training set is the average of individual costs:
    $$ J(\theta) = -\frac{1}{m} \sum_{i=1}^{m} \left[ y^{(i)} \log(h_\theta(x^{(i)})) + (1 - y^{(i)}) \log(1 - h_\theta(x^{(i)})) \right] $$
    \item This function is strictly convex and derivable via Maximum Likelihood Estimation (MLE).
\end{itemize}
\end{frame}

\begin{frame}{Gradient Descent for Logistic Regression}
\begin{itemize}
    \item Our objective is again to find $\min_\theta J(\theta)$ using Gradient Descent.
    \item The update rule is exactly the same as Linear Regression!
    $$ \theta_j := \theta_j - \alpha \frac{\partial}{\partial \theta_j} J(\theta) $$
    $$ \theta_j := \theta_j - \alpha \frac{1}{m} \sum_{i=1}^{m} (h_\theta(x^{(i)}) - y^{(i)}) x_j^{(i)} $$
    \item \textbf{BUT:} It is NOT the same algorithm, because $h_\theta(x)$ is different.
        \begin{itemize}
            \item Linear Reg: $h_\theta(x) = \theta^T x$
            \item Logistic Reg: $h_\theta(x) = \frac{1}{1 + e^{-\theta^T x}}$
        \end{itemize}
\end{itemize}
\end{frame}

\begin{frame}{Advanced Optimization Algorithms}
\begin{itemize}
    \item Gradient Descent works well, but there are more sophisticated, faster algorithms for minimizing $J(\theta)$ that do not require you to manually pick a learning rate $\alpha$.
    \item Examples:
        \begin{itemize}
            \item Conjugate Gradient
            \item BFGS
            \item L-BFGS
        \end{itemize}
    \item Advantages: Often converge much faster than Gradient Descent.
    \item Disadvantages: More complex math. (Luckily, libraries like Scikit-Learn handle them internally).
\end{itemize}
\end{frame}

\begin{frame}{Multi-class Classification: One-vs-Rest (OvR)}
\begin{itemize}
    \item What if we have more than 2 classes? $y \in \{0, 1, 2, \dots, K\}$.
    \item Examples: Email folders (Work, Friends, Family), Weather (Sunny, Cloudy, Rain).
    \item Strategy: \textbf{One-vs-Rest} (or One-vs-All).
    \item We train $K$ separate binary logistic regression classifiers.
    \item Classifier $i$ is trained to predict the probability that the observation belongs to class $i$ vs. all other classes.
    $$ h_\theta^{(i)}(x) = P(y=i | x; \theta) $$
\end{itemize}
\end{frame}

\begin{frame}{Making Predictions in One-vs-Rest}
\begin{itemize}
    \item After training $K$ classifiers, how do we predict the class for a new input $x$?
    \item We pass the new input $x$ to all $K$ classifiers to get $K$ probabilities.
    \item We predict the class $i$ that outputs the highest probability:
    $$ \text{prediction} = \max_i \left( h_\theta^{(i)}(x) \right) $$
\end{itemize}
\end{frame}

\begin{frame}[fragile]{Logistic Regression: Python Demo Setup}
\begin{itemize}
    \item We use Scikit-Learn's \texttt{LogisticRegression} class.
    \item We will use the famous Iris dataset (classifying iris flowers based on petal/sepal dimensions).
\end{itemize}
\begin{minted}[fontsize=\scriptsize]{python}
from sklearn.datasets import load_iris
from sklearn.linear_model import LogisticRegression
from sklearn.model_selection import train_test_split
from sklearn.metrics import accuracy_score, confusion_matrix
\end{minted}
\end{frame}

\begin{frame}[fragile]{Logistic Regression: Python Demo Code}
\begin{minted}[fontsize=\scriptsize]{python}
# 1. Load data
X, y = load_iris(return_X_y=True)

# For a simple binary classification demo, we only keep classes 0 and 1
X_binary = X[y != 2]
y_binary = y[y != 2]

# 2. Split data
X_train, X_test, y_train, y_test = train_test_split(
    X_binary, y_binary, test_size=0.3, random_state=42
)

# 3. Train model (Scikit-learn uses L-BFGS by default)
model = LogisticRegression()
model.fit(X_train, y_train)
\end{minted}
\end{frame}

\begin{frame}[fragile]{Logistic Regression: Python Demo Results}
\begin{minted}[fontsize=\scriptsize]{python}
# 4. Predict and evaluate
y_pred = model.predict(X_test)
acc = accuracy_score(y_test, y_pred)

print(f"Accuracy: {acc * 100:.2f}%")

# Let's see the probabilities for the first 3 test examples
probs = model.predict_proba(X_test[:3])
print("Probabilities (Class 0, Class 1):\n", probs)

# Output might look like:
# Accuracy: 100.00%
# Probabilities:
# [[0.98  0.02]   -> Predicts Class 0
#  [0.05  0.95]   -> Predicts Class 1
#  [0.91  0.09]]  -> Predicts Class 0
\end{minted}
\end{frame}

\begin{frame}{Logistic Regression: Summary (2.2)}
    \begin{itemize}
        \item The go-to baseline algorithm for Binary Classification.
        \item Uses the \textbf{Sigmoid function} to squash linear outputs into probabilities $(0, 1)$.
        \item Uses the \textbf{Cross-Entropy (Log Loss)} cost function to ensure convexity.
        \item Optimized using Gradient Descent or advanced methods (BFGS).
        \item Extends to multi-class problems via the \textbf{One-vs-Rest} strategy.
        \item Prone to overfitting if there are too many polynomial features (can be fixed with Regularization).
    \end{itemize}
\end{frame}

\section{2.3 Support Vector Machines (SVM)}

\begin{frame}{Support Vector Machines: Introduction}
\begin{itemize}
    \item SVM is a powerful and versatile machine learning model, considered one of the most robust supervised learning algorithms.
    \item Originally developed for \textbf{Classification} (SVC).
    \item Can be extended for \textbf{Regression} (SVR) and \textbf{Outlier Detection} (One-class SVM).
    \item Particularly well-suited for classification of complex but small-to-medium-sized datasets.
\end{itemize}
\end{frame}

\begin{frame}{The "Large Margin" Classifier Intuition}
\begin{itemize}
    \item Imagine a binary classification task where data is linearly separable (we can draw a straight line between the two classes).
    \item There are infinite lines that can perfectly separate the training data. Which one is "best"?
    \item Logistic Regression might pick a line that is very close to some points.
    \item SVM takes a different approach: it tries to find the line that is \textbf{farthest away} from the closest training instances.
\end{itemize}
\end{frame}

\begin{frame}{Hyperplane and Margin}
\begin{itemize}
    \item \textbf{Hyperplane:} The decision boundary separating the classes (a line in 2D, a plane in 3D).
    \item \textbf{Margin:} The distance between the hyperplane and the closest data points from either class.
    \item You can think of SVM as fitting the \textbf{widest possible "street"} (represented by parallel lines) between the classes.
    \item The goal is to maximize this margin.
\end{itemize}
\end{frame}

\begin{frame}{Support Vectors}
\begin{itemize}
    \item The instances (data points) that are located exactly on the edge of the "street" (the margin boundaries) are called \textbf{Support Vectors}.
    \item They are critical: they fully define the decision boundary.
    \item If you add or remove other training points (that are off the street), the decision boundary does NOT change.
    \item This makes SVM relatively robust and memory efficient after training.
\end{itemize}
\end{frame}

\begin{frame}{Hard Margin Classification}
\begin{itemize}
    \item If we strictly enforce that all instances must be off the street and on the correct side, this is called \textbf{Hard Margin} classification.
    \item Two major issues with Hard Margin:
    \begin{enumerate}
        \item \textbf{Linear Separability:} It only works if the data is perfectly linearly separable. If classes overlap, it fails.
        \item \textbf{Outlier Sensitivity:} A single outlier can dramatically shrink the margin or make hard margin impossible to find, resulting in poor generalization.
    \end{enumerate}
\end{itemize}
\end{frame}

\begin{frame}{Soft Margin Classification}
\begin{itemize}
    \item To solve the issues of Hard Margin, we need a more flexible model.
    \item \textbf{Soft Margin} allows some instances to end up inside the margin, or even on the wrong side of the hyperplane.
    \item Objective: Find a balance between keeping the margin as large as possible AND limiting \textbf{margin violations}.
\end{itemize}
\end{frame}

\begin{frame}{The Regularization Hyperparameter: $C$}
\begin{itemize}
    \item In Scikit-Learn's SVM classes, this balance is controlled by the hyperparameter $C$.
    \item $C$ acts like an inverse regularization parameter (similar to $1/\lambda$ in Logistic Regression).
    \item \textbf{Small $C$ (More Regularization):} Wider margin, but more margin violations allowed. Good for noisy data (prevents overfitting).
    \item \textbf{Large $C$ (Less Regularization):} Narrower margin, penalizes violations heavily. Forces the model to fit the training data strictly (risk of overfitting).
\end{itemize}
\end{frame}

\begin{frame}{Non-linear Classification}
\begin{itemize}
    \item Real-world data is rarely perfectly linearly separable.
    \item For example, imagine two concentric rings of data points. A straight line cannot separate them.
    \item \textbf{Solution 1:} Similar to polynomial regression, we can add polynomial features (e.g., $x_1^2, x_2^2, x_1x_2$) to transform the 2D data into a 3D (or higher) space where a linear hyperplane \textit{can} separate them.
    \item \textbf{Problem:} Adding many features makes the dataset huge, slowing down computation tremendously.
\end{itemize}
\end{frame}

\begin{frame}{The Kernel Trick (1)}
\begin{itemize}
    \item The \textbf{Kernel Trick} is a brilliant mathematical technique used by SVMs.
    \item It allows the SVM to find a hyperplane in a high-dimensional feature space \textit{without} actually computing the coordinates of the data in that space.
    \item Instead of transforming data $x \rightarrow \phi(x)$ and computing dot products $\phi(a)^T \phi(b)$, a Kernel function $K(a,b)$ computes the result directly.
\end{itemize}
\end{frame}

\begin{frame}{The Kernel Trick (2) - Benefits}
\begin{itemize}
    \item \textbf{Computational Efficiency:} It saves enormous amounts of memory and processing power.
    \item \textbf{Infinite Dimensions:} It allows SVMs to operate in an implicit feature space with infinite dimensions (like the RBF kernel), which would be impossible to compute explicitly.
    \item Common Kernels: Linear, Polynomial, Gaussian RBF, Sigmoid.
\end{itemize}
\end{frame}

\begin{frame}{Gaussian RBF Kernel (1)}
\begin{itemize}
    \item The Gaussian \textbf{Radial Basis Function (RBF)} is the most popular kernel for non-linear SVMs.
    \item It measures the "similarity" or distance between two instances.
    \item Formula: $K(a, b) = \exp(-\gamma ||a - b||^2)$
    \item It essentially places a bell-shaped curve over each training instance (a "landmark").
\end{itemize}
\end{frame}

\begin{frame}{Gaussian RBF Kernel (2) - The $\gamma$ hyperparameter}
\begin{itemize}
    \item \textbf{$\gamma$ (gamma)} defines how far the influence of a single training example reaches.
    \item \textbf{Small $\gamma$:} The bell curve is wide. Instances have a far reach. The decision boundary becomes smoother (more regularized).
    \item \textbf{Large $\gamma$:} The bell curve is narrow. Instances only influence very close neighbors. The decision boundary becomes very irregular, tightly wrapping around data points (risk of overfitting).
\end{itemize}
\end{frame}

\begin{frame}{Tuning SVM: $C$ and $\gamma$}
\begin{itemize}
    \item When using an SVM with an RBF kernel, tuning $C$ and $\gamma$ simultaneously is critical.
    \item Usually done via \textbf{Grid Search} with Cross-Validation.
    \item If the model overfits: Decrease $C$ and/or decrease $\gamma$.
    \item If the model underfits: Increase $C$ and/or increase $\gamma$.
    \item Feature Scaling (like StandardScaler) is \textbf{mandatory} for SVMs, as they are based on distance calculations (margin).
\end{itemize}
\end{frame}

\begin{frame}{SVM Mathematical Formulation (Primal Problem)}
\begin{itemize}
    \item \textit{(Optional/Advanced)} The mathematical objective of a linear SVM (Hard Margin).
    \item Let hyperplane be $w^T x + b = 0$. Margin size is $\frac{2}{||w||}$.
    \item Maximize margin $\rightarrow$ Minimize $\frac{1}{2} ||w||^2$.
    \item Constraint: $y^{(i)}(w^T x^{(i)} + b) \ge 1$ for all $i=1 \dots m$.
    \item This is a \textbf{Convex Quadratic Programming (QP)} problem, solved efficiently by specialized solvers (like LibSVM).
\end{itemize}
\end{frame}

\begin{frame}{Hinge Loss}
\begin{itemize}
    \item An alternative way to view SVM is as optimizing a loss function using Gradient Descent.
    \item SVM uses the \textbf{Hinge Loss} function: $\max(0, 1 - t)$
    \item Where $t = y^{(i)}(\theta^T x^{(i)})$.
    \item If the instance is correctly classified and outside the margin ($t \ge 1$), loss is 0.
    \item If it's inside the margin or misclassified ($t < 1$), the loss increases linearly.
\end{itemize}
\end{frame}

\begin{frame}{Support Vector Regression (SVR)}
\begin{itemize}
    \item SVM can be adapted for regression (predicting continuous numbers).
    \item \textbf{SVR Trick:} Reverse the objective! Instead of trying to fit the largest possible street \textit{between} two classes while limiting margin violations...
    \item SVR tries to fit \textbf{as many instances as possible ON the street} while limiting margin violations (instances \textit{off} the street).
    \item The width of the street is controlled by a hyperparameter $\epsilon$ (epsilon).
\end{itemize}
\end{frame}

\begin{frame}[fragile]{SVM: Python Demo Setup}
\begin{itemize}
    \item We will use Scikit-Learn to classify handwritten digits (0-9).
    \item Note the importance of scaling features before applying SVM.
\end{itemize}
\begin{minted}[fontsize=\scriptsize]{python}
from sklearn import datasets
from sklearn.svm import SVC
from sklearn.model_selection import train_test_split
from sklearn.preprocessing import StandardScaler
from sklearn.metrics import accuracy_score, classification_report
\end{minted}
\end{frame}

\begin{frame}[fragile]{SVM: Python Demo Code}
\begin{minted}[fontsize=\scriptsize]{python}
# 1. Load digits dataset (8x8 images of numbers)
digits = datasets.load_digits()
X, y = digits.data, digits.target

# 2. Split data
X_train, X_test, y_train, y_test = train_test_split(
    X, y, test_size=0.3, random_state=42
)

# 3. Mandatory Feature Scaling for SVM
scaler = StandardScaler()
X_train_scaled = scaler.fit_transform(X_train)
X_test_scaled = scaler.transform(X_test)

# 4. Train SVM (using Gaussian RBF kernel)
svm_clf = SVC(kernel='rbf', C=1.0, gamma='scale')
svm_clf.fit(X_train_scaled, y_train)
\end{minted}
\end{frame}

\begin{frame}[fragile]{SVM: Python Demo Results}
\begin{minted}[fontsize=\scriptsize]{python}
# 5. Predict and evaluate
y_pred = svm_clf.predict(X_test_scaled)
acc = accuracy_score(y_test, y_pred)

print(f"SVM (RBF Kernel) Accuracy: {acc * 100:.2f}%\n")

# Detailed classification report
print(classification_report(y_test, y_pred))

# Example Output:
# SVM (RBF Kernel) Accuracy: 98.15%
# (Followed by precision, recall, f1-score for digits 0-9)
\end{minted}
\end{frame}

\begin{frame}{Support Vector Machines: Summary (2.3)}
    \begin{itemize}
        \item \textbf{Core Idea:} Maximize the margin between classes (Large Margin Classifier).
        \item \textbf{Support Vectors:} Only a subset of data points dictates the decision boundary, making it memory efficient.
        \item \textbf{Flexibility:} Can handle strict separation (Hard Margin) or noisy data (Soft Margin via hyperparameter $C$).
        \item \textbf{Kernel Trick:} Efficiently handles complex, non-linear data mapping to high dimensions (e.g., RBF kernel via $\gamma$).
        \item \textbf{Requirement:} Features must be scaled (Standardization) before training.
    \end{itemize}
\end{frame}

\end{document}
